\subsection{Sparse ProxQP for lifted interval constraints}\label{sec:method-proxqp}
ProxQP solves a convex quadratic program using a primal-dual proximal
augmented Lagrangian (Bambade et al., \href{https://simple-robotics.github.io/publications/proxqp/}{primary source}).
The adapter forms the same separable slack or group-peak epigraph lift
used for the other sparse comparators. It partitions a lifted row into
an equality row only when its finite lower and upper bounds coincide.
The canonical backend problem is
\begin{equation}\label{eq:proxqp-interval-form}
 \min_z \ \tfrac12z^\top P_Lz+q_L^\top z,\qquad
 Ez=e,\quad Cz\in B=\prod_i[l_i,u_i],
 \qquad P_L\succeq0.
\end{equation}
The lifted variables include the costed auxiliaries, so splitting the
hard interval rows does not alter which soft violations incur cost.
A missing interval side is represented by an infinite endpoint.

One way to understand the inequality update is to introduce \(w=Cz\)
with \(w\in B\). Let \(v_k\) be the current signed inequality multiplier
and \(\mu>0\) its proximal multiplier parameter. For a fixed \(z\), the
augmented quadratic in \(w\) is minimized by the interval projection,
giving the equivalent signed multiplier update
\begin{equation}\label{eq:proxqp-dual-update}
 \widehat w=\Pi_B(Cz+\mu v_k),\qquad
 v^+=v_k+\mu^{-1}(Cz-\widehat w).
\end{equation}
For the free equality multiplier \(y_k\), the analogous update is
\(y^+=y_k+\mu_e^{-1}(Ez-e)\) with \(\mu_e>0\).
Together with a positive primal proximal term
\(\rho\lVert z-z_k\rVert_2^2/2\), these updates yield a regularized,
piecewise quadratic primal subproblem. The projection changes its
local derivative when a bound becomes active. ProxQP uses the
primal-dual merit system and an active-region Newton step with a line
search to solve the proximal subproblems inexactly. The published
convergence discussion depends on a primal-dual solution with zero
duality gap and controlled inner errors. The adapter does not extend
those conditions to indefinite or otherwise unsupported inputs.

\begin{algorithm}[ht]
\caption{Sparse ProxQP and signed recovery}\label{alg:proxqp-route}
\begin{algorithmic}[1]
\State Assemble the lifted Hessian, equality block and remaining interval block.
\State Initialize the sparse solver without an initial guess.
\State Solve the proximal primal-dual subproblems until the native termination test.
\If{status is \texttt{PROXQP\_SOLVED}}
 \State Reorder equality and signed interval multipliers to their original lifted rows.
 \State Check lifted feasibility, stationarity and penalty recovery in physical coordinates.
\Else
 \State Retain the native termination without an accepted candidate.
\EndIf
\end{algorithmic}
\end{algorithm}

ProxSuite 0.7.2 receives sparse matrices, no initial guess, absolute
and relative residual and gap tolerances of \(10^{-7}\), an iteration
limit of \(200\,000\), and gap checking. Its interval multiplier has
upper-positive and lower-negative sign. Bound activity alone is not a
basis for supplying a missing multiplier. Sparse matrix assembly is
linear in stored entries, while numerical cost is governed by the
factorizations, their fill, and the inner and outer steps. Fresh
instances avoid sharing an inferred active set or numerical
factorization across samples.
